\section{Spectral caps and comparison with spherical spins}\label{cap:static}
The cap route first compares partition functions and projected means with a
Gaussian measure conditioned on a sphere. The comparison is uniform over
macroscopic balls of edge fields. Its role is to control the values and slopes
needed when differentiating a Gaussian observation convolution.

\subsection{Observation and scale conventions}
For a symmetric matrix $A$ and field $h$, write
\begin{equation}\label{cap:partition-definition}
 Z_A(h)=2^{-n}\sum_{x\in\{-1,1\}^n}
       \exp\{\tfrac12 x^\top A x+h^\top x\},\qquad
 \phi_A(h)=\log Z_A(h).
\end{equation}
We write $Z_p=Z_{J_p}$ once $J_p$ has been defined, and use $Z_A^{\rm sp}$ for
the analogous integral against uniform probability on the sphere of radius
$\sqrt n$. Gradients and Hessians of $\phi_A$ are the Gibbs mean and
covariance.

Here are the conventions for all cap sections. Every tolerance described as
arbitrarily small is a fixed positive number chosen before $n\to\infty$. We
write $\eta_{\rm err}$ for such a tolerance, which can be decreased at later
finite stages of the parameter selection; $\eta_{\rm small}$ has the same
meaning when a separate structural slack is needed. Neither symbol denotes a
rate in $n$. Once the scale ratios, field ranges, and error tolerances have
been fixed, the upper scale is chosen small enough, and only then is $n$ taken
large. Uniform $o(1)$ terms refer to this order of limits. High probability,
unless otherwise specified, only means probability tending to one. Variables
called fields are in energy units. A Gaussian observation with precision \(B\)
is the field increment \(Bx+N(0,B)\). Its effect on posterior log weights is
to add the field and subtract \(B\) from the interaction matrix. One can use
increasing precisions and independent noise increments (including Brownian
observations). The logarithms of partition functions before an increment are
obtained from those after it by log heat convolution. Additive constants in
potentials don't affect gradients, heat-tilted transition probabilities or
Hessians. We use Euclidean matrix and vector norms; empirical averages have
brackets, and sometimes we use \(|v|_n=|v|/\sqrt n\).

The edge coordinates are eigenvectors with gap \(d_i=2-\lambda_i\), in order
of increasing gap. We condition in spectral-coordinate arguments on the
eigenvalues, leaving the frame Haar. The observation scheme is given by the
spectral functions
\begin{equation}
\label{cap:edge-definitions}\quad P_p=1_{\{d \le p^2\}},\qquad k_p=np^3,\qquad M_p=\sqrt{np},
 \qquad B_p=(p^2-d)_+,\qquad J_p=J-B_p .
\end{equation}
Write \(H_p=\operatorname{ran}P_p\) and \(d_p=\dim H_p=\operatorname{Tr}P_p\);
the spectral estimates below give \(d_p\asymp k_p\). We use fields in
\(H_p\), and occasionally write \(P_p\) for this space. Field length
units there are \(p^2 M_p\); free-energy and projected-spin length units are
\(k_p,M_p\). The parameter \(p\) is small. Conditionally on \(J_p\) and the
eigenvalues, refinements determining smaller caps use a fresh Haar rotation
within \(P_p\).

\subsection{Spectral estimates at the bottom scale}

The resolvent event of Lemma~\ref{prof:lem:spectrum} covers the entire
cap range. Indeed, if $np_{\min}^3\asymp\sqrt{\log n}$, then
\[
 \frac{p_{\min}}{d_0}\asymp
 \frac{(\log n)^{1/6}}{\log\log(n+1000)}\longrightarrow\infty.
\]
Thus the lemma applies at $d=\sqrt u$ for every $u\ge c p^2$, where
$p\ge p_{\min}$ and $c>0$ is fixed. In particular, on one ordinary
disorder event of probability tending to one,
\[
 d_1\ge-o(p^2),\qquad \|J\|\le3,
 \qquad \#\{i:d_i\le u\}\le Cnu^{3/2}\quad(p^2\le u\le4).
\]
The count follows by applying the lemma on a dyadic grid in $\sqrt u$
and using monotonicity between grid points. The ordinary empirical
spectral measure also converges to the semicircle law.

Write $m_J(u)=n^{-1}\operatorname{Tr}(d+u)^{-1}$. The same lemma and
$b_{\rm sc}(2+u)=1-\sqrt u+O(u)$ give, uniformly for small
$u\ge cp^2$,
\[
 m_J(u)=b_{\rm sc}(2+u)
       +o\!\left(\frac{\sqrt u}{\sqrt{nu^{3/2}}}\right)
 =1-\bigl(1+o(1)+O(\sqrt u)\bigr)\sqrt u.
\]
Here $nu^{3/2}\ge c^{3/2}np_{\min}^3\to\infty$.

We will also use the precise edge counting measure. Put
\[
 \nu_p=\frac1{k_p}\sum_i\delta_{d_i/p^2}.
\]
For fixed $v,w>0$, its transform difference is exactly
\[
 \int\left(\frac1{x+v}-\frac1{x+w}\right)d\nu_p(x)
 =\frac{m_J(vp^2)-m_J(wp^2)}p
 =\sqrt w-\sqrt v+o(1)+O(p).
\]
The limiting expression equals the same integral against
$\pi^{-1}\sqrt x\,\mathbf1_{\{x>0\}}dx$. The count bound above controls
the tail of the integrand, which is $O(x^{-2})$, by $O(R^{-1/2})$ beyond
$R$. Since the negative support lies in $[-o(1),0]$, the measures
$(1+x)^{-2}\nu_p$ are tight. Uniqueness of their Stieltjes transforms
therefore yields
\[
 \nu_p\ \longrightarrow\ \pi^{-1}\sqrt x\,
                    \mathbf1_{\{x>0\}}dx
 \quad\hbox{on bounded intervals}.
\]
All these statements are uniform in the cap range in the convention
that $n\to\infty$ after the fixed upper cutoff on $p$ is chosen, and
then that cutoff may be decreased. For example, a failure of uniformity
would give a sequence of cap scales on which the uniform resolvent
estimate and the same tightness argument give the displayed limit,
a contradiction.

The spectral event is imposed once for all scales. Later failures of
size $e^{-a k_p}$ on a geometric cap list are also summable:
\[
 \sum_{j\ge0}e^{-a\gamma^{-3j}k_{\min}}=o(1),
 \qquad k_{\min}\asymp\sqrt{\log n}.
\]

\subsection{Spherical saddles, norm densities, and tails}
We compare several times with spins uniform on the Euclidean sphere of radius
\(\sqrt n\), with superscript \({\rm sp}\) for this change. Write
\(\mu_{p,h}\) and \(\mu_{p,h}^{\rm sp}\) for the cube and spherical Gibbs
probabilities normalized by \(Z_p(h)\) and \(Z_p^{\rm sp}(h)\).
At the cap \(p\) with field \(h\in H_p\) in any fixed bounded range of field
units, write
\begin{equation}
\label{cap:saddle}\quad R_{p,\alpha}=(\max(d,p^2)+\alpha)^{-1},\qquad
 a_{p,\alpha}=R_{p,\alpha}h,\qquad
 n={\rm Tr}\,R_{p,\alpha}+\|a_{p,\alpha}\|^2 .
\end{equation}
This fixes the spectral parameter \(\alpha=\alpha_p(h)\). At this solution
we abbreviate \(R_p=R_{p,\alpha}\) and \(a_p=a_{p,\alpha}\). For any admissible
\(\alpha\), let \(f_{p,\alpha,h}^{\rm norm}\) denote the density of
\(\|N(R_{p,\alpha}h,R_{p,\alpha})\|^2\). The spherical posterior is
\(N(a_p,R_p)\) conditioned on squared norm \(n\). One has
\(-1+c\le\alpha/p^2\le C\) by the spectral facts. Conversely lengths for
\(\alpha=0\) satisfy, with \(b=4/(3\pi)\),
\begin{equation}
\label{cap:trace-constants}\quad n-{\rm Tr}\,R_p=(b+o(1))np,\qquad {\rm Tr}\,R_p^2=(2b+o(1))n/p .
\end{equation}
Changing \(\alpha\) by a small fraction of \(p^2\) perturbs these by relative
small amounts. For bounded arguments in scaled units the saddle and its
inverse (squared field length as function of parameter, where nonnegative) are
uniformly Lipschitz in scaled variables.

We record density-estimate details used below. Changing quadratic and linear
tilts that retain a comparable spectral band, at fixed \(\alpha\), uses
ordinary Gaussian mgfs times the
ratio of squared-norm densities at \(n\), by polar coordinates. At the cap
saddle the density at \(n\) is comparable to \(\sqrt{p/n}\). Indeed the
norm-square variance is of order \(n/p\); in standard deviation units the
characteristic function converges locally to \(e^{-t^2/2}\) after centering
(Taylor expansion, max Gaussian coordinate variance divided by norm-square
standard deviation \(O(1/\sqrt{k_p})\)). Absolute values are bounded by \((1+c
t^2/k_p)^{-c k_p}\) from top components alone, also in the noncentral case.
This proves the assertion by Fourier inversion. The sphere restriction is
imposed by conditioning a Gaussian squared norm. To transfer an
exponential-moment estimate through this conditioning, we need a lower density
at the original saddle and an upper density after the additional tilt, both on
the same scale. The characteristic function below supplies these two bounds.
Its central limit expansion gives the density near the mean, while an entire
spectral band supplies an integrable bound away from the origin.

In diagonal coordinates, for $Y_i=a_i+\sqrt{r_i}g_i$ the formula used is
\begin{equation}\label{cap:norm-characteristic}
 \left|\E e^{it\sum_iY_i^2}\right|
 =\prod_i(1+4t^2r_i^2)^{-1/4}
   \exp\left\{-\sum_i\frac{2t^2r_i a_i^2}{1+4t^2r_i^2}\right\}.
\end{equation}
A band of order $k_p$ coordinates with $r_i\asymp p^{-2}$ supplies the
integrable majorant, and the remaining factors have absolute value at most
one. The variance formula $\operatorname{Var}(\sum_iY_i^2)=2\Tr
R_p^2+4a_p^\top R_pa_p$ sets the scale $n/p$. This proves the centered lower
density bound. It also gives the noncentral upper bound for a linear tilt,
and for a quadratic tilt whenever the new covariance retains a band with the
same size and comparable eigenvalues.

The same argument applies if an extra scalar multiple of \(p^2 P_p\) is added
to the interaction keeping comparable gaps at the saddle. Upper bounds suffice
without the exact trace constraint or centering.

We will use the corresponding bound for a fixed number \(m\) of replicas.
Let independent spectral rows \(Y_i\in\mathbb R^m\) have laws
\(N(a_i,C_i)\) with \(C_i\succ0\), and put
\(\mathcal Q_j=\sum_iY_{ij}^2\). For \(T_t=\operatorname{diag}(t_1,\ldots,t_m)\),
\(H_i=C_i^{1/2}T_tC_i^{1/2}\), and \(\beta_i=C_i^{-1/2}a_i\), the Gaussian
quadratic formula gives
\[
 \left|\E e^{iY_i^\top T_tY_i}\right|
 =\det(I+4H_i^2)^{-1/4}
   e^{-2\beta_i^\top H_i^2(I+4H_i^2)^{-1}\beta_i}
 \le \det(I+4H_i^2)^{-1/4}.
\]
Suppose \(L>2\) rows satisfy \(c s^{-1}I_m\preceq C_i\preceq C s^{-1}I_m\).
The singular values of \(H_i\) are at least \(c s^{-1}\) times those of
\(T_t\), even when \(T_t\) has both signs. Independence therefore yields
\[
 \left|\E e^{it\cdot\mathcal Q}\right|
 \le\prod_{j=1}^m(1+c t_j^2/s^2)^{-L/4}.
\]
Each one-dimensional integral is at most \(Cs/\sqrt L\). Fourier inversion
then proves
\begin{equation}\label{cap:joint-norm-density}
 \sup_q f_{\mathcal Q}(q)\le C_m(s/\sqrt L)^m
 \le C_m(s^{1/4}/\sqrt n)^m
 \quad\text{if }L\asymp ns^{3/2}.
\end{equation}
Independent replicas at their original cap saddles have joint density at
\(n\mathbf1\) at least \(c_m(p/n)^{m/2}\). Thus an additional Gaussian
tilt that preserves the displayed band costs at most
\(C_m(\sqrt s/p)^{m/2}\) in the norm-density ratio.

The density in this last comparison must be the density under the
normalized tilted law. Indeed, for \(0\le F\le1\) with \(\E F>0\),
\[
 \E[F\mid\mathcal Q=n\mathbf1]
 =\E F\,\frac{f_{\mathcal Q}^{F}(n\mathbf1)}
                  {f_{\mathcal Q}(n\mathbf1)},
 \qquad d\mathbb P_F=\frac{F\,d\mathbb P}{\E F},
\]
where \(f_{\mathcal Q}^{F}\) is the norm density under \(\mathbb P_F\);
at \(n\mathbf1\) we use the representatives given by polar coarea.
For a tail \(F=\mathbf1_{\{\mathcal T\ge b\}}\), exponential Markov gives
the useful polar-conditioning form
\begin{equation}\label{cap:conditioned-tail}
 \Pr(\mathcal T\ge b\mid\mathcal Q=n\mathbf1)
 \le e^{-\lambda b}\E e^{\lambda\mathcal T}
       \frac{f_{\mathcal Q}^{(\lambda)}(n\mathbf1)}
            {f_{\mathcal Q}(n\mathbf1)},\qquad
 d\mathbb P^{(\lambda)}
 =\frac{e^{\lambda\mathcal T}d\mathbb P}{\E e^{\lambda\mathcal T}} .
\end{equation}
Here the conditional quantities are defined by polar disintegration.
For a rowwise linear or quadratic \(\mathcal T\), the tilted law remains
Gaussian: a factor \(e^{Y_i^\top A_iY_i/2+v_i^\top Y_i}\) changes the
covariance to \((C_i^{-1}-A_i)^{-1}\). A strict positive precision margin
that keeps a full band comparable to \(s^{-1}I_m\) permits
\eqref{cap:joint-norm-density} for \(f_{\mathcal Q}^{(\lambda)}\).
For a general projected test, its spherical failure probability will instead
be an explicit hypothesis in the test-transfer argument below.

We also record the exact linear moment bound at a cap saddle. Polar
conditioning with the parameter frozen at \(\alpha=\alpha_p(h)\) gives
\begin{equation}\label{cap:frozen-mgf}
 \log\E_{\mu_{p,h}^{\rm sp}}e^{v^\top X}
 =a_p^\top v+\tfrac12v^\top R_pv+
   \log\frac{f_{p,\alpha,h+v}^{\rm norm}(n)}
                 {f_{p,\alpha,h}^{\rm norm}(n)}
 \le a_p^\top v+\tfrac12v^\top R_pv+C_{\rm dens}.
\end{equation}
The denominator uses the centered lower bound and the numerator the
noncentral upper bound with the same covariance. Thus \(C_{\rm dens}\)
is independent of \(n,p\) on the fixed bounded ranges used here.

Here are some tail estimates in this convention. For a fixed finite number of
spherical replicas at a bounded cap tilt as above (also allowing the
just-mentioned scalar modification), scaled projection norms \(\|P_p x\|/M_p\)
above large \(v\) cost \(\exp(-c k_p v^6)\), up to a decrease in \(c\); pair
overlaps above \(r\ge C p\), \(r\) small, cost \(\exp(-c n r^3)\). The same
norm tail applies for the completely uncapped zero-field model. For detail,
change the Gaussian parameter upward by \(s=\epsilon r^2\) for overlaps (take
\(r=v^2p\le1\) for projection norms), \(\epsilon\) small. The cost in log
normalizer for this is \(O(n s^{3/2})\) for \(s\gg p^2\), by integrating the
trace deficit; bounded cap fields add at most \(O(k_p)\). At uncapped zero
field one can start at parameter \(\rho p^2\), fixed small \(\rho>0\): the
cost in the lower partition bound via its norm density is at most \(C k_p\) in
addition to the log of the density scale \(\sqrt{p/n}\) (fill the \(O(np)\)
deficit with one top Gaussian square, the other squares concentrating by
Chebyshev).

In the changed Gaussian the mean contributes negligibly and the desired tails
cost \(c n r^3\epsilon\), by exponential Markov with parameter \(c's\) in the
quadratic form (trace-square \(O(n/\sqrt s)\)).

For this exponential tilt the replica precision on a band of order
\(ns^{3/2}\) rows remains comparable to \(sI\), if the Markov parameter
\(c's\) is a sufficiently small multiple of \(s\). Equations
\eqref{cap:conditioned-tail} and \eqref{cap:joint-norm-density} therefore
bound the joint norm-density ratio for \(m\) replicas by
\(C_{m,\epsilon}(r/p)^{m/2}\); for a pair it is \(C_\epsilon r/p\).
This introduces no separate factor in \(n\) when \(r\asymp p\).
For \(r\ge p\) its logarithm is absorbed by the tail exponent because
\(nr^3=k_p(r/p)^3\) and \(k_p\to\infty\). More precisely, for fixed \(a>0\),
\[
 C_{m,\epsilon}(r/p)^{m/2}e^{-a nr^3}
 \le C_{m,\epsilon,a}e^{-a nr^3/2}
\]
uniformly once \(k_p\) is large. Taking \(\epsilon\) small and then \(C\)
large gives the assertions. For overlaps at cap saddles with
zero field one also has below order \(p\) cost \(\exp(O(1)-c n p r^2)\) by the
same density/mgf argument at the saddle. For angular-moment comparisons at
overlaps bounded away from zero we give details separately below.

\subsection{Uniform comparison of values and means}
Comparing the Ising and spherical models by rotation has precedents in
Comets~\cite{Comets1996Spherical}; critical free-energy and overlap
comparisons also appear in Du--Huang~\cite{DuHuang2026CriticalFluctuations,
DuHuang2026CriticalOverlap}. The statement below includes the field and
posterior-test uniformity required by the covariance argument.
\begin{proposition}[Static cap comparison]\label{cap:static-comparison}
Fix a field radius $R<\infty$, an error tolerance $\eta>0$, and a geometric
ratio in $(0,1)$. There is a sufficiently small upper scale $p_0>0$ such that,
with probability tending to one, simultaneously over a geometric list of caps
$p\le p_0$ with $np^3\ge c\sqrt{\log n}$ and all
$h\in P_p\mathbb R^n$ satisfying $\|h\|\le R p^2 M_p$, one has
\begin{equation}\label{cap:static-values}
 |\log Z_p(h)-\log Z_p^{\rm sp}(h)|\le \eta k_p,\qquad
 \|P_p m_p(h)-a_p(h)\|\le \eta M_p,
 \qquad m_p(h)=\nabla\log Z_p(h).
\end{equation}
The projected norm tails proved above hold, with a decreased positive
exponential constant, under these cube posteriors and under the uncapped
zero-field cube posterior. The statement remains valid for any fixed finite collection of auxiliary
geometric cap lists. The zero-field concentration and density estimates
used in the proof also allow bounded scalar modifications of the cap
interaction whose spherical saddle gaps remain comparable to $p^2$.
\end{proposition}
\begin{proof}
It suffices to prove the value comparison on a slightly larger prescribed
ball, with an error \(\varepsilon_v k_p\) chosen below the requested value
error. To see the mean implication, let \(u\in H_p\) be a unit vector and
choose \(t>0\) so that \(h\pm tu\) stay in the larger ball. The two value
bounds and \eqref{cap:frozen-mgf} give
\[
 \log\E_{\mu_{p,h}}e^{\pm t u^\top(X-a_p)}
 \le 2\varepsilon_v k_p+\tfrac12t^2u^\top R_pu+C_{\rm dens}.
\]
Jensen's inequality in the two signs yields
\[
 |\langle u,P_pm_p(h)-a_p\rangle|
 \le \tfrac t2u^\top R_pu+
       \frac{2\varepsilon_v k_p+C_{\rm dens}}t.
\]
Set \(t=\theta p^2M_p\), and take \(u\) along the projected discrepancy.
The saddle separation gives \(p^2\|R_p|_{H_p}\|\le C\), while
\(p^2M_p^2=k_p\). Hence
\[
 \frac{\|P_pm_p(h)-a_p\|}{M_p}
 \le C\theta+\frac{2\varepsilon_v+C_{\rm dens}/k_p}{\theta}.
\]
Choose \(\theta\) from the desired mean tolerance, then
\(\varepsilon_v\ll\theta^2\), and finally \(k_p\) large. The larger ball
provides the same fixed margin at boundary fields. We now prove the value
comparison and the stated test transfers.

\paragraph{Replica Gram comparison.}
 Haar first moments agree exactly for cube and sphere. For \(l\)-th moments, \(l\) fixed, the angular average once the replica Gram matrix is given is again the same. At off-diagonal normalized overlaps bounded by small \(\rho\) the cube Gram law is bounded by the spherical law, with factor
\begin{equation}
                       \exp(C_l+C_l n\rho^4),
 \label{cap:gram-comparison}
\end{equation}
using comparison cells of radius \(O_l(1/n)\). Indeed the off-diagonal
products of the signs in a row form a bounded lattice vector in fixed
dimension \(l(l-1)/2\), with identity covariance. Exponential tilting gives
near the origin a local upper probability bound \(C_l n^{-l(l-1)/4}\) times
the Cramér exponential. The prefactor follows by Fourier inversion (off
neighborhoods of the finite dual-lattice points in the torus one has strict
decay, uniformly in small tilts). Its entropy rate agrees with
\(-\frac12\log\det Q\), \(Q\) the Gram matrix with diagonal one, up to and
including degree three: the only nonzero third moments are triangles, giving
the cubic term minus the sum of triangle monomials. The spherical density
proportional to \((\det Q)^{(n-l-1)/2}\), by successive projection of unit
vectors, has the reciprocal-width prefactor. Cells of sidelength comparable to
\(1/n\) can be taken disjoint. Cholesky coordinates on the angular integral
have bounded derivatives near the identity so change of the exponential
weights on the cells costs \(O_l(1)\) in logarithm. Inequalities on projected
samples used as restrictions there can correspondingly be thickened by
\(O_l(1/n)\) in unit-sphere norm.

For two replicas and \(\rho\le n^{-0.27}\), the binomial masses equal
spherical masses on overlap cells of width \(2/n\) to relative \(o(1)\). At
zero field the angular integrand is even increasing in absolute overlap (the
sum/difference orthogonal coordinates have weights with coefficients
\(1+r,1-r\)); thus these cell comparisons can then use monotonicity instead of
a constant-factor loss.

\paragraph{Zero field and overlaps bounded away from zero.}
Fix the eigenvalues and let \(\E_O\) denote Haar averaging of the eigenframe.
At uncapped zero field, and also for a small zero-field cap with the bounded
scalar modifications above and comparable saddle separation, we have the
following formulas. Here \(A\) denotes the uncapped interaction \(J\) or
the indicated capped interaction.
\[
 \E_O Z_A(0)=Z_A^{\rm sp}(0),\qquad
 \frac{\E_O Z_A(0)^2}{(Z_A^{\rm sp}(0))^2}=1+o(1).
\]
For the second assertion, outside \(n^{-0.27}\) up to a small constant,
\eqref{cap:gram-comparison} costs less than the spherical tail gain on dyadic
bands. For the uncapped sphere use the preceding overlap argument with
lower-normalizer scale \(p\) taken minimal. At overlaps bounded away from
zero the net exponential is strictly negative even with binomial weights.
Here is one verification. In sum and difference orthogonal coordinates,
upper exponents are bounded by the spherical free energies at \(1+|r|\) and
\(1-|r|\); restriction to the orthogonal hyperplane does not change the
bound by interlacing. In the limit at zero cutoff the spherical upper bound
\(F(t)\), exact at \(1\), has \(F'(t)=t/2\) for \(0<t\le1\) and
\(F'(t)=1-1/(2t)\) above \(1\). This follows by Gaussian change of density
with precision \(t(w-J)\), giving per-site bound
\((tw-1-\log t-n^{-1}\log\det(w-J))/2\), optimized at \(w=t+1/t\) below
one and at \(w=2+0^+\) above. The per-site density correction for the upper
bound at strictly positive separation is \(o(1)\), also in the hyperplane;
equality at \(1\) follows by the deficit-filling lower bound. Hence
\(F(1+r)+F(1-r)-2F(1)<r^2/2\), whereas the binomial cost is at least
\(r^2/2\). Continuity handles endpoints, and a small operator-norm change
of the matrices preserves these strict bounds.

All errors giving convergence of the moment ratio are summable over our
geometric list (Stirling errors negative powers of \(n\); tails for \(|r|\ge
n^{-0.27}\) exponentially small in a power of \(n\)). This proves zero-field
concentration and transfers projection tails, by the first moment and Markov.

\paragraph{Positive moments and uniform upper bounds.}
The zero-field comparison supplies a denominator and removes pairs with
macroscopic overlap. At a nonzero field, we first keep the replicas in the
small-overlap region and compare their joint Gram law. The resulting
high-moment estimates have two uses: they control the upper value of the
partition function, and they transfer a posterior test whose spherical
failure has a positive exponential rate. The latter use requires an
additive estimate on discarded replicas after the denominator has been
bounded below.

For bounded field lengths, high positive moments \emph{restricted to small
overlaps among all replicas} have normalized exponents as small as desired in
units \(k_p\), by \eqref{cap:gram-comparison} and the spherical tail at
overlaps above \(C p\); the relative error \(O_l(np^4)\) here suffices. The
restriction discards a negligible relative contribution with high probability
in the quenched model, uniformly: pairs at overlap above fixed small cutoff
have exponentially small probability at cap zero field by the preceding
argument, and the field costs only \(O(n p^{5/2})\). One may take the upper
scale small after fixing \(l\).

Here is the normalized form of the test transfer used later. Fix the
eigenvalues and a field \(u\) in the spectral copy of \(H_p\); in physical
coordinates \(J_p(O)=O J_p^{\rm diag}O^\top\) and \(h_O=Ou\).
Let \(E_u\) be one of the projected-configuration events described by the
Gram-cell comparison, with fixed geometric slack. Write \(E_u^+\) for its
specified thickening: whenever a configuration satisfies \(E_u\), all
configurations in the same Gram cell satisfy \(E_u^+\). Define the tested,
unnormalized sum
\[
 Z_{p,E,O}(u)=2^{-n}\sum_x
 e^{x^\top J_p(O)x/2+h_O^\top x}\,
       \mathbf1_{E_u}(O^\top x).
\]
Its ratio to \(Z_{p,O}(u)=Z_p(h_O)\) is the posterior tested probability.
Assume that the spherical probability of \(E_u^+\) is at most \(e^{-c k_p}\).
For \(l\) independent copies, let
\(\mathcal Z_{E,l}^{\rm small}(O,u)\) be the sum defining
\(Z_{p,E,O}(u)^l\), restricted to pair overlaps below the fixed small
cutoff. The Gram comparison and the spherical tail give
\begin{equation}\label{cap:tested-moment}
 \frac{\E_O\mathcal Z_{E,l}^{\rm small}(O,u)}
      {(Z_p^{\rm sp}(h_O))^l}
 \le \exp\{(-lc/2+\eta_l)k_p+o_n(k_p)\}.
\end{equation}
For fixed \(l\), the exponent loss \(\eta_l>0\) can be made as small as
prescribed by decreasing the upper scale; this includes the
\(O_l(np^4)\) Gram error. The remaining \(o_n(k_p)\) is taken afterward.
The spherical denominator depends only on the spectral field \(u\).
On overlap bands larger than \(Cp\), Cauchy--Schwarz divides the spherical
test cost by two and the remaining spherical overlap cost pays
\eqref{cap:gram-comparison}. This is the source of the factor \(1/2\).

Let \(b_l(O,u)\) be the posterior \(l\)-replica probability that some pair
exceeds the fixed overlap cutoff. Directly from the decomposition of the
replica sum,
\[
 \left(\frac{Z_{p,E,O}(u)}{Z_{p,O}(u)}\right)^l
 \le \frac{\mathcal Z_{E,l}^{\rm small}(O,u)}{Z_{p,O}(u)^l}
       +b_l(O,u).
\]
On the quenched high-probability zero-field event, the pair estimate above,
followed by the bounded-field change of density, makes \(b_l\) exponentially
small in \(n\), uniformly after the upper scale is chosen small for fixed
\(l\). Once the value lower bound
proved below gives \(Z_{p,O}(u)\ge e^{-\varepsilon_v k_p}Z_p^{\rm sp}(h_O)\),
Markov's inequality applied to \eqref{cap:tested-moment} shows that
\begin{align}\label{cap:tested-transfer}
 &\Pr_O\left\{
 \frac{Z_{p,E,O}(u)}{Z_{p,O}(u)}>2e^{-c'k_p},\
 Z_{p,O}(u)\ge e^{-\varepsilon_v k_p}Z_p^{\rm sp}(h_O),\
 b_l(O,u)\le e^{-2lc'k_p}\right\}\nonumber\\
 &\hspace{25mm}\le
 \exp\{[-l(c/2-c'-\varepsilon_v)+\eta_l]k_p+o_n(k_p)\}.
\end{align}
Choose \(c'<c/2\) and \(\varepsilon_v<c/2-c'\). Increasing the fixed
replica number \(l\), and only then decreasing the upper scale, gives any
prescribed fixed Haar failure rate. The separately bounded event for \(b_l\)
is added after normalization, as the displayed decomposition requires.
The untested positive moments give the upper value bounds.

The reusable test statement takes the spherical rate as a hypothesis.
Linear and quadratic projection tests, including bounded spectral
transformations, supply that rate through \eqref{cap:conditioned-tail}
when their exponential tilt retains the required precision margin.
When a later concrete test also uses Gaussian observation noise, each
replica receives an independent noise mark. Its application must check the
same cellwise inclusion with those marks fixed; a small spherical probability
alone does not extend the statement to an arbitrary bounded test.

Upper value bounds extend uniformly by convexity from a fine constant mesh on
a larger ball. The zero-field bound (and evenness) then gives a Lipschitz
constant in scaled units on inner balls. Posterior tests with slack and
positive-rate tails can likewise use nets, by restricting the change of
density to bounded projection lengths first (the removed large-length tails
follow from the zero-field version and absorption). We only use nets of size
\(\exp(O(k_p))\).

\paragraph{Angular lower bounds and gradients.}
Upper moments alone do not provide a lower value at every field direction.
We first obtain a set of directions on each length sphere where the value
is large and the projected gradient is close to its prediction. The scalar
modification of the cap turns a zero-field partition comparison into an
average over such a sphere. A second-moment estimate then controls how
small this set of directions can be.

First, on each fixed length sphere in field space, up to arbitrarily small log
losses in scaled units its angular partition average has the predicted lower
bound with high probability. To see this at nonzero length add a scalar
interaction on \(P_p\) to \(J_p\) so that, at the parameter \(\alpha\)
predicted for the given length, the eigenvalue of \(R_p\) on \(P_p\) increases
by \(\|a_p\|^2/{\rm Tr}P_p\). This restores the trace to \(n\), with
comparable separation. Decouple the addition by a centered Gaussian. Under the
integrated spherical weight the field length thus generated concentrates
exponentially at the specified length (the unconditioned generated field using
the modified saddle is Gaussian, with mean zero and expected squared length
precisely \(\|h\|^2\); the norm constraint has at most polynomial cost in
\(k_p\), by density bounds, or by a small quadratic tilt with the bounds
above). By first moment and the zero-field two-moment argument at the modified
cap the same shell has nearly all the integrated cube weight. Narrowing the
shell to any fixed tolerance and using the scaled Lipschitz bounds gives the
claim.

Angular second moments of \(Z/Z^{\rm sp}\) are at most \(e^{\eta_{\rm
err}\,k_p}\) with high probability by the small-overlap comparison. Moreover
the two-replica numerator cost is \(e^{-c_\tau k_p}\) for
\begin{equation}
 |(P_p x-a_p)^t(P_p x'-a_p)|>\tau M_p^2
\end{equation}
as follows. In the two independent canonical Gaussians, the exponential
factor \(e^{\pm\theta p^2(P_pX-a_p)^\top(P_pX'-a_p)}\) changes each
two-replica precision by an off-diagonal block of size \(\theta p^2\) on
\(H_p\). For sufficiently small \(\theta\) the cap band remains comparable
to \(p^2I_2\), and the Gaussian log moment is at most \(C\theta^2 k_p\).
Equation~\eqref{cap:conditioned-tail} then gives exponent
\(-\theta\tau k_p+C\theta^2k_p+O(1)\); choosing \(\theta\) as a small
multiple of \(\tau\) gives the stated spherical rate. The Gram comparison
transfers it to the angular numerator. Polynomial weights in the scaled
projections are handled by the projected norm tails.

Thus there is a set of angular measure at least \(e^{-\eta k_p}\), for
arbitrarily small prescribed \(\eta\), where the lower value holds up to error
\(\eta k_p\) and gradients differ from \(a_p\) in projection by any prescribed
small fraction of \(M_p\). Indeed discard the points with low value and those
with large normalized mean discrepancy in the angular first moment (use the
two-replica bound and Cauchy-Schwarz), then use the angular second moment.
Losses in these arguments can be chosen in order after the tolerances. Length
grids suffice. At these directions, the value and gradient estimates define
supporting planes. Evaluating those planes at another field gives a lower
bound there. The Gaussian moment-generating-function bound and the bounded
saddle range show that the extra scaled loss is a small multiple of the
distance plus a constant multiple of its square.

\paragraph{Refinement from a coarser cap.}
The remaining task is to reach every child field from the large-value sets
just constructed. Freeze the child saddle parameter and use its canonical
Gaussian observation as a reference bridge to a coarser cap. Under this
reference law the parent length is predictable. Fresh Haar refinement makes
the reference field close to the useful parent directions, and their
supporting planes convert that distance estimate into a lower bound for
the child value. The gain must pay for a net at the child scale; this is why
the additional scale ratio is chosen small.

For this comparison it is useful to retain the exact normalization of the
frozen Gaussian quadratic. For any admissible \(\alpha\), define
\begin{equation}\label{cap:frozen-potential}
 \begin{split}
 C_{n,\alpha}
 &=\log\!\left(2^{n/2}\Gamma(n/2)n^{1-n/2}\right)
       +\tfrac n2(2+\alpha),\\
 q_{u,\alpha}^{\rm fr}(h)
 &=C_{n,\alpha}+\tfrac12\log\det R_{u,\alpha}
                   +\tfrac12h^\top R_{u,\alpha}h .
 \end{split}
\end{equation}
The determinant is on the ambient space. Coarea in squared radius gives,
for every such \(\alpha\),
\begin{equation}\label{cap:sphere-frozen-density}
 \phi_u^{\rm sp}(h)
 =q_{u,\alpha}^{\rm fr}(h)+\log f_{u,\alpha,h}^{\rm norm}(n).
\end{equation}
This identity also follows by writing the Gaussian density on
\(\|x\|^2=n\): its factor independent of \(x\) is
\((2\pi)^{-n/2}(\det R_{u,\alpha})^{-1/2}
e^{-n(2+\alpha)/2-h^\top R_{u,\alpha}h/2}\), and the squared-radius
coarea factor is \(\pi^{n/2}n^{n/2-1}/\Gamma(n/2)\).

To obtain a uniform lower bound at \(c\), write \(c=rp\) for a very small
fixed ratio \(r\), using an additional coarser cap list. Put
\(\Delta=B_p-B_c\), freeze \(\alpha=\alpha_c(h_c)\), and abbreviate
\(R_u=R_{u,\alpha}\), \(a_c=R_ch_c\). Since
\(R_c^{-1}=R_p^{-1}-\Delta\), Gaussian integration gives
\[
 q_{c,\alpha}^{\rm fr}(h_c)
 =\log\int e^{q_{p,\alpha}^{\rm fr}(h_c+z)}N(0,\Delta)(dz).
\]
The corresponding probability law for the parent field is
\begin{equation}\label{cap:static-canonical-bridge}
 \begin{split}
 \mathsf G_{c,p}^{h_c}(dy)
 &=e^{q_{p,\alpha}^{\rm fr}(y)-q_{c,\alpha}^{\rm fr}(h_c)}
                  N(h_c,\Delta)(dy)\\
 &=N(h_c+\Delta a_c,\ \Delta+\Delta R_c\Delta)(dy),\qquad
 R_pY\sim N(a_c,R_c-R_p).
 \end{split}
\end{equation}
Equivalently, sample \(X\sim N(a_c,R_c)\) without norm conditioning and set
\(Y=h_c+\Delta X+N(0,\Delta)\). Then
\(X\mid Y=y\sim N(R_py,R_p)\). Conditioning this joint law on
\(\|X\|^2=n\) gives the spherical transition and the exact density ratio
\begin{equation}\label{cap:spherical-bridge-ratio}
 \frac{d\mathsf T_{c,p}^{\rm sp}}{d\mathsf G_{c,p}^{h_c}}(y)
 =\frac{f_{p,\alpha,y}^{\rm norm}(n)}
        {f_{c,\alpha,h_c}^{\rm norm}(n)} .
\end{equation}

The numerator here is generally evaluated away from the parent's own
saddle. Let \(\widehat\alpha=\alpha_p(y)\). From
\eqref{cap:sphere-frozen-density},
\[
 \log\frac{f_{p,\alpha,y}^{\rm norm}(n)}
              {f_{p,\widehat\alpha,y}^{\rm norm}(n)}
 =q_{p,\widehat\alpha}^{\rm fr}(y)-q_{p,\alpha}^{\rm fr}(y).
\]
The first \(\alpha\)-derivative of \(q_{p,\alpha}^{\rm fr}(y)\) is
\(\tfrac12(n-\Tr R_{p,\alpha}-\|R_{p,\alpha}y\|^2)\), which vanishes at
\(\widehat\alpha\). Its second derivative is
\(\tfrac14\operatorname{Var}_{p,\alpha,y}(\|X\|^2)\). Thus, with
\(d_\alpha=\alpha-\widehat\alpha\),
\begin{equation}\label{cap:off-saddle-density}
 \log\frac{f_{p,\alpha,y}^{\rm norm}(n)}
              {f_{p,\widehat\alpha,y}^{\rm norm}(n)}
 =-\frac{d_\alpha^2}{4}\int_0^1(1-t)
  \operatorname{Var}_{p,\widehat\alpha+t d_\alpha,y}(\|X\|^2)\,dt .
\end{equation}
On bounded ranges with comparable gaps this variance is \(O(n/p)\).
Consequently a shell \(|\widehat\alpha-\alpha|\le\tau p^2\) costs at most
\(C\tau^2k_p\) in the logarithm. The centered parent and child densities
have ratio of order \(\sqrt{p/c}\), so their logarithmic cost is \(O_r(1)\).
This is a statement about the ratio: each individual density contains the
factor \(n^{-1/2}\).

The true heat convolution, rewritten using \(\mathsf G=\mathsf G_{c,p}^{h_c}\),
is
\[
 \phi_c(h_c)-q_{c,\alpha}^{\rm fr}(h_c)
 =\log\E_{\mathsf G}e^{\phi_p(Y)-q_{p,\alpha}^{\rm fr}(Y)} .
\]
Jensen and \eqref{cap:sphere-frozen-density} therefore give
\begin{equation}\label{cap:frozen-jensen}
 \phi_c(h_c)-\phi_c^{\rm sp}(h_c)
 \ge \E_{\mathsf G}\bigl[\phi_p(Y)-\phi_p^{\rm sp}(Y)\bigr]
   +\E_{\mathsf G}\log
       \frac{f_{p,\alpha,Y}^{\rm norm}(n)}
            {f_{c,\alpha,h_c}^{\rm norm}(n)} .
\end{equation}
We must bound both expectations from below, including the reference fields
outside any fixed shell. We first prove that the normalization term satisfies
\begin{equation}\label{cap:averaged-frozen-density}
 -C_{R,r}\le
 \E_{\mathsf G}\log
       \frac{f_{p,\alpha,Y}^{\rm norm}(n)}
            {f_{c,\alpha,h_c}^{\rm norm}(n)}
 \le0 .
\end{equation}
Here $R$ is the prescribed child field radius. Since $r$ is fixed and
$k_c=r^3k_p\to\infty$, this loss is $o(k_c)$ even at the bottom scale.

\paragraph{Averaged normalization loss.}
The child saddle bounds give $-1+\kappa_R\le\alpha/c^2\le C_R$,
with constants depending on $R$ before $r$ is chosen. Taking $r$ small
makes $|\alpha|/p^2$ small. Put
\[
 v=(p^2+\alpha)^{-1},\qquad D=n-\Tr R_p.
\]
Then $d_p\asymp k_p$, $v\asymp p^{-2}$, $D\asymp np$, and
$\Tr R_p^2\asymp n/p$. For the deficit, changing the parameter from
zero to $\alpha$ changes the trace by at most
$C|\alpha|n/p=O_R(r^2np)$, while \eqref{cap:trace-constants} gives
the positive deficit $(b+o(1))np$ at zero. The same estimates hold in
a sufficiently small fixed $p^2$-neighborhood of $\alpha$.

We need a density lower bound valid for every $y\in H_p$. Split the
squared norm of $N(R_py,R_p)$ as $S_a+S_b$, where the independent
active and bulk parts have laws
\[
 S_a=\|N(vy,vI_{H_p})\|^2,\qquad
 S_b=\|N(0,R_p|_{H_p^\perp})\|^2.
\]
Write $D_b=n-\E S_b=D+d_pv\asymp np$. Since
$\Var S_b\le Cn/p$, Chebyshev gives
\[
 \Pr\{|S_b-\E S_b|\le D_b/2\}\ge1-C/k_p\ge\tfrac12.
\]
On this event $t=n-S_b\in[D_b/2,3D_b/2]$. Spherical integration
in the active space, followed by Jensen with a uniform unit vector
$\omega\in H_p$, gives
\[
 f_{S_a}(t)
 =e^{-v\|y\|^2/2}f_{v\chi_{d_p}^2}(t)
       \E_\omega e^{\sqrt t\,y^\top\omega}
 \ge e^{-v\|y\|^2/2}f_{v\chi_{d_p}^2}(t).
\]
Here $t/(d_pv)$ lies in a fixed compact subset of $(0,\infty)$.
Stirling's bound for the chi-square density therefore gives
$f_{v\chi_{d_p}^2}(t)\ge C^{-1}(v\sqrt{d_p})^{-1}e^{-Ck_p}$.
Convolving over the bulk event proves
\begin{equation}\label{cap:global-frozen-density}
 f_{p,\alpha,y}^{\rm norm}(n)
 \ge C^{-1}\sqrt{p/n}\,
       \exp\left\{-Ck_p\left(1+
              \frac{\|y\|^2}{p^4M_p^2}\right)\right\}
 \qquad(y\in H_p).
\end{equation}
The child density is comparable to $\sqrt{c/n}$, and the noncentral
upper bound from \eqref{cap:norm-characteristic} bounds the parent
density by $C\sqrt{p/n}$ for every $y$. Their common $n^{-1/2}$ factor
thus cancels before taking logarithms. In particular the logarithmic
ratio in \eqref{cap:averaged-frozen-density} is integrable under the
Gaussian bridge.

For its sharper average bound, set $T=R_c-R_p$ and
$W=\|R_pY\|^2$. The bridge and the child saddle equation give
\[
 \E_{\mathsf G}W=\|a_c\|^2+\Tr T=D,\qquad
 \Var_{\mathsf G}W=2\Tr T^2+4a_c^\top Ta_c\le C_Rn/c.
\]
Indeed $\Tr R_c^2\le C_Rn/c$, $\|R_c\|\le C_Rc^{-2}$, and
$\|a_c\|^2\le C_Rnc$. Fix a sufficiently small $\delta>0$ and put
$\mathcal C=\{|W-D|\le\delta np\}$. The Gaussian quadratic moment
formula bounds $\log\E e^{t(W-D)}$ by $C_Rt^2n/c$ for
$|t|\le c_Rc^2$. Exponential Markov in both signs consequently gives
$\mathsf G(\mathcal C^c)\le C_{R,r,\delta}e^{-a_{R,r,\delta}k_p}$.

On $\mathcal C$, monotonicity of the saddle equation places
$\widehat\alpha=\alpha_p(Y)$ in the comparable-gap neighborhood above:
the derivative of $\Tr R_{p,s}+\|R_{p,s}Y\|^2-n$ has negative
magnitude at least $c_2n/p$ there. Integrating that derivative to each
endpoint first locates the root, and then gives
\[
 |\widehat\alpha-\alpha|\le C(p/n)|W-D|.
\]
Along the intervening parameter segment, the norm-square variance is
at most $Cn/p$. The centered density bounds and
\eqref{cap:off-saddle-density} imply on $\mathcal C$ that
\begin{align*}
 \log\frac{f_{p,\alpha,Y}^{\rm norm}(n)}
              {f_{c,\alpha,h_c}^{\rm norm}(n)}
 &\ge-C_{R,r}-C(n/p)|\widehat\alpha-\alpha|^2,\\
 \E_{\mathsf G}\!\left[
  \mathbf1_{\mathcal C}(n/p)|\widehat\alpha-\alpha|^2\right]
 &\le C(p/n)\Var W\le C_Rp/c.
\end{align*}
For the complement, $\|Y\|^2/(p^4M_p^2)\asymp W/(np)$ has bounded
second moment. Cauchy--Schwarz and the preceding exponential tail show
that the negative contribution allowed by
\eqref{cap:global-frozen-density} is at most
$C_{R,r,\delta}k_pe^{-a_{R,r,\delta}k_p/2}=o(1)$.
This proves the lower bound in \eqref{cap:averaged-frozen-density}.
Integrability and \eqref{cap:spherical-bridge-ratio} identify its
expectation with $-D_{\rm KL}(\mathsf G\Vert\mathsf T_{c,p}^{\rm sp})$,
which proves the upper bound.

\paragraph{Averaging the parent supporting planes.}
On a finite length grid near the typical parent length, let $\mathcal U$
be the closed set of fields $z$ satisfying
\[
 \phi_p(z)-\phi_p^{\rm sp}(z)\ge-\varepsilon_ak_p,
 \qquad \|P_pm_p(z)-a_p(z)\|\le\varepsilon_gM_p.
\]
Choose $\varepsilon_a,\varepsilon_g$ using the preceding angular
argument, so that these sets have angular measure at least
$e^{-\eta k_p}$ on each grid sphere, with $\eta$ as small as required.
The set $\mathcal U$ uses only parent data and the prescribed radii,
so it is fixed under the fresh refinement. These anchors lie in a fixed
bounded parent field range. Convexity gives
a lower supporting plane for $\phi_p$ at $z$. On the spherical side,
\eqref{cap:frozen-mgf}, frozen at the own saddle of $z$, gives
\[
 \phi_p^{\rm sp}(y)\le\phi_p^{\rm sp}(z)
       +a_p(z)^\top(y-z)
       +\tfrac12(y-z)^\top R_{p,\alpha_p(z)}(y-z)+C_{\rm dens}.
\]
This holds for every $y\in H_p$: the denominator density is centered at
$z$, and the numerator has the same covariance, so its noncentral upper
bound is independent of $y$. Subtracting the two inequalities yields the
global estimate
\begin{equation}\label{cap:global-anchor-support}
 \phi_p(y)-\phi_p^{\rm sp}(y)
 \ge-\varepsilon_ak_p-\varepsilon_gM_p\|y-z\|
       -Cp^{-2}\|y-z\|^2-C_{\rm dens}
 \qquad(y\in H_p,\ z\in\mathcal U).
\end{equation}
Put $d(y)=\operatorname{dist}(y,\mathcal U)/(p^2M_p)$ and
$D_2=\E_{\mathsf G}d(Y)^2$. Choosing a nearest anchor and averaging
\eqref{cap:global-anchor-support} gives
\begin{equation}\label{cap:averaged-anchor-support}
 \E_{\mathsf G}[\phi_p(Y)-\phi_p^{\rm sp}(Y)]
 \ge-k_p\{\varepsilon_a+\varepsilon_g\sqrt{D_2}+CD_2\}
       -C_{\rm dens}.
\end{equation}
In particular this average includes all Gaussian length tails.

To bound this distance, put $\ell_*=(p^2+\alpha)\sqrt D$.
The child saddle $\alpha$ depends only on the eigenvalues and
$\|h_c\|$, since $h_c\in H_c$; thus $\ell_*$ is independent of the
fresh orientation. The preceding variance estimate gives
\[
 \E_{\mathsf G}\left[
   \frac{(\|Y\|-\ell_*)^2}{p^4M_p^2}\right]
 \le\frac{C\Var W}{Dnp}=O_{R,r}(k_p^{-1}).
\]
Let the grid mesh be at most $\xi p^2M_p$. Under Haar averaging the
direction of $Y$ is uniform conditional on its length. Choose a grid
radius within $\xi p^2M_p$ of $\ell_*$ and apply spherical
concentration to its anchor set of angular measure at least
$e^{-\eta k_p}$. The radial bound and the triangle inequality give
\[
 \E_O D_2\le C_R(\eta+\xi^2)+O_{R,r}(k_p^{-1}),
\]
where $O$ denotes the fresh refinement. Hence the Haar average of
$\sqrt{D_2}$ is as small as prescribed. Its Frobenius Lipschitz constant is
at most $C\sqrt r+O_r(k_p^{-1/2})$. Indeed the reference field in parent
units has mean norm $O(\sqrt r)$ and covariance norm $O_r(1/k_p)$;
the Gaussian $L^2$ norm of its change between two rotations bounds the
change of its $L^2$ distance to the fixed set $\mathcal U$.
Haar Lipschitz concentration therefore gives
$D_2\le\epsilon r^3+o(1)$ except with probability
$\exp(-c_3\epsilon r^2k_p)$ (proper rotations suffice).
The constants before choosing $r$ depend only on the prescribed child
field range. First choose $\epsilon$ from the requested child value
error $\varepsilon_v$, then choose $r$ sufficiently small that the
failure rate pays for the fixed fine child net, whose size is
$\exp(O(k_c))=\exp(O(r^3k_p))$. Next choose
$\varepsilon_a,\varepsilon_g,\eta,\xi$ so that the Haar mean of
$\sqrt{D_2}$ is smaller than $\tfrac12\sqrt{\epsilon r^3}$ and the
coefficient $\varepsilon_a+\varepsilon_g\sqrt{D_2}+CD_2$ is a
sufficiently small multiple of $\varepsilon_vr^3$ on the retained event.
Then choose the upper cap scale sufficiently small for the preceding
angular estimates, and finally take $n$ large with all parameters fixed.
Equations
\eqref{cap:frozen-jensen}, \eqref{cap:averaged-frozen-density}, and
\eqref{cap:averaged-anchor-support} give the child lower bound throughout
the net. Extend by the scaled Lipschitz bound.
All finitely many grids and extra lists, with positive-rate failures per
$k_p$ or the summable zero-field errors above, retain high probability.
This proves \eqref{cap:static-values}.

\end{proof}
